\section{Background and Related Work}\label{sec:background}

Research in non-prehensile manipulation dates back to the 90s \cite{lynch1999dynamic, ruggiero2018nonprehensile, toussaint2018differentiable}, with the vast majority of prior work leveraging heuristics or analytical models based on simplifying assumptions. Related work has focused on skills such as throwing \cite{zeng2020tossingbot}, poking \cite{Huang-1997-14452}, and pushing \cite{woodruff2017planning,RSS2020Changkyu,mason1986mechanics}.
More recently, pushing manipulation has received increased attention due to availability of large-scale datasets and the inherent controllability of the skill \cite{bauza2019omnipush}. Pushing operates under the quasistatic assumption (i.e., the inertial effects of robot--object and object--environment interaction are ignored) to reduce modeling complexity, thereby limiting robot velocities and accelerations. On the contrary, poking must plan around the non-negligible effects of inertial forces: the object continues moving after robot--object contact is broken, thus allowing for faster planar manipulation of objects. Past work in pushing manipulation incorporates simulation and analytical models in a motion planning loop to get the next feasible state \cite{zito2012two}\cite{zhou2019pushing}. Additional contributions in push modeling involve combining object state estimation with affordance prediction from image data to determine contact points for achieving the optimal push \cite{kloss2020accurate} and creating a deep recurrent neural network model to model push outcomes for a variety of objects \cite{li2018push}. However, both approaches use a greedy planner operating in obstacle-free environments. In this work, we propose a sampling-based kinodynamic framework that is capable of planning dynamic collision-free paths in the object configuration space. 

Past work on sampling-based planning techniques for non-prehensile manipulation involve kinodynamic approaches to rearrangement planning and whole-arm manipulation. This work is constrained  by the quasistatic assumption, thereby limiting object manipulation speed through the use of non-dynamic primitives \cite{king2015nonprehensile}. However, when operating in a dynamic regime, lack of kinematic modeling of the robot and lack of contact modeling between the robot arm and environmental objects may result in unusable actions explored during the planning phase \cite{haustein2015kinodynamic}. These approaches also employ an open-loop paradigm and do not take robot, sensing, and model uncertainty into account. Moreover, open-loop kinodynamic planners that consider uncertainty in the planning process may generate conservative plans that do not exploit robot and object dynamics to their full extent \cite{8207585}. Our approach uses a simulation engine that considers robot kinematics and contact while planning and operates in a closed-loop manner, replanning if the resultant pose is outside a certain threshold.
% \ap{Additionally, a key difference to prior work is the explicit consideration of manipulator constraints and the extension \textsl{PokeRRT*} that is used to produce shorter solutions to leverage poking's crucial ability to cover large distances quickly.}
% Additionally, neither approach leverages RRT*, thereby not exploiting the long action paths that are crucial for taking full advantage of dynamic primitive-based planning. -- but here is work on kinodynamic RRT* so probably shouldn't mention this

In all, evidence from prior work suggests that \textsl{poking}---sometimes referred to as ``releasing'' or ``impulsive manipulation''---is a relatively unexplored primitive. Analytical models for poking dynamics are restricted to rotationally symmetric objects or situations where pusher--object contact can be geometrically modeled \cite{Huang-1997-14452, zhu2005frictional}. The need for specialized impulse-delivery apparatus to achieve poking is explored in \cite{Huang-1997-14452}. This makes planar manipulation of objects cumbersome due to the manual relocation of the apparatus required. In this work, we use a standard open-chain robot arm equipped with an electric parallel gripper to deliver impulses using joint space velocity control.

In order to generate impulse-based action paths that obey the robot's kinematic constraints and produce feasible object motion, we use a PyBullet \cite{coumans2019} simulation environment as the forward physics model in the planning loop.
% Since methods for skill modeling are inextricably tied to the research on NPM planning detailed above, as a dynamical model is a prerequisite of any NPM system.
% Specifically, we use PyBullet \cite{coumans2019} as the forward physics model in our planner. 
While simulation may not accurately capture real-world dynamics \cite{collins2019quantifying}, the closed-loop nature of our planning framework compensates for this approach. This effective use of simulation captures essential characteristics of robot and object dynamics at faster than realtime speeds and in a safe manner without inducing significant wear-and-tear caused by dynamic trajectory execution on a real robot.
